%=============================================================================!
% 11.E  EXERCISES                                                             !
%=============================================================================!
\unnumbered{Exercises}
\label{sec:cs-exercises}

The exercises run from questions of understanding, through calculations,
to short derivations marked `show that'. Full solutions are in
\hyperref[sec:cs-solutions]{`Solutions to the exercises'} at the end of
the book, and Notebook~11 checks every numerical answer.

\begin{exercise}[freedoms of methanol]
\label{ex:cs-freedoms}
Methanol, CH$_3$OH, has six atoms and four bonds to hydrogen. How many
freedoms has one molecule with those four bonds held, and how many have
100 such molecules once the drift of their centre of mass is removed?
\end{exercise}

\begin{exercise}[one bond solved exactly]
\label{ex:cs-one-bond}
Atoms $i$ and $j$, of masses 2 and \SI{5}{\amu}, are joined by a bond held
at \SI{1.5}{\angstrom}, with $\vb{r}_{ij} = (1.5, 0, 0)$\,\AA\ at the start
of a step. The provisional bond is $\tilde{\vb{r}}_{ij} = (1.6, 0.3,
0)$\,\AA. Solve \cref{eq:cs-quadratic} for $\eta$, choose the root, and
find how far each atom moves.
\end{exercise}

\begin{exercise}[Newton's method for a square root]
\label{ex:cs-newton}
Apply Newton's method to $f(y) = y^2 - 3$ from $y_0 = 2$, and give $y_1$,
$y_2$, $y_3$ and their errors against $\sqrt3$.
\end{exercise}

\begin{exercise}[the error squares]
\label{ex:cs-newton-error}
Show that for $f(y) = y^2 - a$, with $a > 0$, Newton's method gives
$y_{n+1} = \frac12(y_n + a/y_n)$, and that $y_{n+1} - \sqrt a = (y_n -
\sqrt a)^2/(2y_n)$ exactly.
\end{exercise}

\begin{exercise}[SHAKE keeps the centre of mass]
\label{ex:cs-centre}
Show that a SHAKE correction of one bond, for any $\eta$, leaves the
centre of mass of its two atoms where it was.
\end{exercise}

\begin{exercise}[what the velocity correction removes]
\label{ex:cs-kinetic}
For two atoms, show that their kinetic energy is $\frac12M\lvert\vb{V}
\rvert^2 + \frac12m_{\mathrm r}\lvert\vb{v}_{ij}\rvert^2$, with $M$ and
$\vb{V}$ their total mass and the velocity of their centre of mass. Hence
show that the velocity correction of RATTLE leaves the total momentum
unchanged and lowers the kinetic energy by $\frac12m_{\mathrm r}
(\vb{v}_{ij}\cdot\vb{e})^2$, with $\vb{e}$ the unit vector along the
bond.
\end{exercise}

\begin{exercise}[correcting an O--H pair]
\label{ex:cs-oh}
An oxygen atom of mass \SI{16}{\amu} and a hydrogen of mass \SI{1}{\amu}
are joined by a bond held at \SI{0.96}{\angstrom}, with $\vb{r}_{\mathrm H}
- \vb{r}_{\mathrm O} = (0.96, 0, 0)$\,\AA, and velocities $\vb{v}_{\mathrm O} = (0.001, 0, 0)$ and $\vb{v}_{\mathrm H} =
(0.02, 0.01, 0)$\,\AA/fs. Find $\zeta$, the corrected velocities and the
kinetic energy removed, in amu\,\AA$^2$/fs$^2$.
\end{exercise}

\begin{exercise}[a larger box]
\label{ex:cs-count}
How many freedoms have 216 rigid water molecules once the drift is
removed, and what kinetic energy, in \si{\electronvolt}, gives each of them
\SI{12.93}{\milli\electronvolt}?
\end{exercise}

\begin{exercise}[one inner step]
\label{ex:cs-respa-one}
Show that RESPA with one inner step is velocity Verlet with the force
$\vb{F}^{\mathrm f} + \vb{F}^{\mathrm s}$.
\end{exercise}

\begin{exercise}[the cost of a nanosecond]
\label{ex:cs-cost}
With \SI{24}{\milli\second} per call of the forces between molecules and
\SI{0.04}{\milli\second} per call of the springs, how long do velocity
Verlet at \SI{0.41}{\femto\second} and RESPA with outer steps of
\SI{0.72}{\femto\second} and three inner steps take for a nanosecond of
flexible water?
\end{exercise}

\begin{exercise}[the square of a matrix]
\label{ex:cs-matrix}
Check that $\vb{M}^2 = (\operatorname{tr}\vb{M})\vb{M} - (\det\vb{M})
\vb{I}$ for the matrix with rows $(1, 2)$ and $(3, 4)$.
\end{exercise}

\begin{exercise}[the second band]
\label{ex:cs-band}
Show that, for small $b = (\omega_{\mathrm s}/\omega_{\mathrm f})^2$, the
trace of \cref{eq:cs-trace} exceeds 2 for $\theta$ from $2\pi(1 - b)$ to
$2\pi$, a band of width $b\,\mathcal{T}_{\mathrm f}$ in the outer step.
\end{exercise}

\begin{exercise}[predicted and measured bands]
\label{ex:cs-bands}
For $\omega_{\mathrm s} = 0.3\,\omega_{\mathrm f}$, give the predicted
edges of the two bands of \cref{fig:cs-resonance}, in units of
$\mathcal{T}_{\mathrm f}$, and compare them with 0.459 to 0.500 and 0.919
to 1.000 from the full trace.
\end{exercise}

\begin{exercise}[a finer inner step]
\label{ex:cs-inner}
With outer steps of \SI{2}{\femto\second}, RESPA on flexible water gives
an energy fluctuation of 8.8\% with inner steps of
\SI{0.25}{\femto\second} and 10.4\% with \SI{0.125}{\femto\second}.
Explain why the finer inner step does not help.
\end{exercise}

\begin{exercise}[hydrogen doubled]
\label{ex:cs-hmr-double}
The hydrogen atoms of water are made twice as heavy, \SI{2.016}{\amu},
at the oxygen's expense. Find the new mass of the oxygen, the O--H
reduced mass, and the factor by which the stretch slows.
\end{exercise}

\begin{exercise}[where the centre of mass goes]
\label{ex:cs-hmr-centre}
How far from the oxygen, along the line bisecting the H--O--H angle, is
the centre of mass of a TIP3P molecule with ordinary masses, and with the
hydrogens three times as heavy?
\end{exercise}

\begin{exercise}[what repartitioning changes]
\label{ex:cs-hmr-what}
Which of these change when hydrogen masses are repartitioned: the mean
potential energy, the distribution of distances between oxygen atoms, the
rate at which molecules diffuse, the frequencies of the vibrations?
\end{exercise}

\begin{exercise}[a nanosecond of drift]
\label{ex:cs-drift}
The faulty start of \cref{sec:cs-trajectory} gives 64 water molecules
of \SI{18.015}{\amu} a total momentum $\vb{P}$ with components 0.2817,
0.2936 and 0.4033\,amu\,\AA/fs. How far does the centre of mass move in a nanosecond,
and how much kinetic energy does the drift carry?
\end{exercise}

\begin{exercise}[a step without its unit]
\label{ex:cs-ase-step}
ASE measures time in units of \SI{10.18}{\femto\second}. A step written
as 2 without converting it to that unit is how long, and what is $\omega
\dt$ for the O--H stretch of period \SI{8.88}{\femto\second}? Is velocity
Verlet stable?
\end{exercise}

\begin{exercise}[naming the fault]
\label{ex:cs-diagnose}
A run of rigid water at \SI{2}{\femto\second} has a mean total energy that
falls steadily by \SI{0.5}{\electronvolt\per\pico\second}, by the same
amount at \SI{1}{\femto\second}; its bond errors are $10^{-10}$\,\AA\ and
its centre of mass is still. Which fault of \cref{tab:cs-faults} fits, and
what test of the model would confirm it?
\end{exercise}
